% STATUS: written 2026-11-05
% Sources read in full: Kolchin 1948 (Ann. of Math. 49, 1--42, secs. 17--20);
% Hubbard-Lundell 2011 (Amer. Math. Monthly 118, 245--261, sec. 4);
% Cassidy-Singer 2007 (arxiv version, sec. 2); Michel--Srinivasan 2026
% (pp. 1--3, added 2026-11-06).
% Sources: Kolchin1948PicardVessiot, HubbardLundell2011, CassidySinger2007,
% MichelSrinivasan2026.
\section{Picard--Vessiot and differential Galois theory}\label{sec:differential-algebra:differential-galois}

Section~\ref{sec:differential-algebra:ritt-kolchin} attached algebraic
varieties to systems of differential equations. The Galois question is
different: pick one homogeneous linear differential equation over a
differential field, and ask which group describes the algebraic relations
among its solutions \cite[pp.~1--6]{Kolchin1948PicardVessiot}. Picard and
Vessiot saw what the answer should look like
\cite[pp.~1--2]{Kolchin1948PicardVessiot}. Kolchin built it in 1948
\cite[pp.~1--6]{Kolchin1948PicardVessiot}.

Kolchin works over an ordinary differential field $F$ of characteristic
zero whose field of constants $C$ is algebraically closed
\cite[p.~29]{Kolchin1948PicardVessiot}. To a homogeneous linear
differential polynomial $L(y) = y^{(n)} + p_1 y^{(n-1)} + \dots + p_n y$
with $p_i\in F$ he attaches an extension of $F$ generated by a full
system of solutions, with no new constant
\cite[p.~29]{Kolchin1948PicardVessiot}.

\begin{definition}[{\cite[$\S$17, p.~29]{Kolchin1948PicardVessiot}}]\label{def:diffalg:pv-extension}
A differential extension field $G$ of $F$ is a \emph{Picard--Vessiot
extension} of $F$ for the homogeneous linear differential polynomial
$L(y)$ if
\begin{enumerate}
\item $L(y)$ has a fundamental system of solutions $t_1, \dots, t_n$ with
  $G = F\langle t_1, \dots, t_n\rangle$;
\item the field of constants of $G$ is $C$.
\end{enumerate}
\end{definition}

The group of interest is the group of differential automorphisms of $G$
over $F$ \cite[pp.~29--31]{Kolchin1948PicardVessiot}. Acting on a
fundamental system of solutions, such an automorphism sends each $t_i$ to
a $C$-linear combination of the $t_j$ \cite[p.~29]{Kolchin1948PicardVessiot}.
So the group embeds in $\mathrm{GL}_n(C)$
\cite[p.~29]{Kolchin1948PicardVessiot}. Kolchin shows that its image is
closed in the Zariski topology.

\begin{theorem}[{\cite[$\S$17, Theorem, p.~31]{Kolchin1948PicardVessiot}}]\label{thm:diffalg:pv-group-algebraic}
The group of differential automorphisms of a Picard--Vessiot extension
$G$ over $F$ is isomorphic to an algebraic matric group with a manifold
in $C^{n^2}$.
\end{theorem}

The group therefore inherits the structure theory of algebraic matric
groups of the first chapter of the same paper: normal chains, solvability,
Lie reducibility to triangular form
\cite[pp.~16--24]{Kolchin1948PicardVessiot}. Hubbard and Lundell state
the same fact for the numerical setting
\cite[p.~250]{HubbardLundell2011}.

\begin{theorem}[{\cite[Theorem~4.3, p.~250]{HubbardLundell2011}}]\label{thm:diffalg:dgal-algebraic}
The differential Galois group of a linear differential operator is an
algebraic subgroup of the general linear group of its solution space:
a subset defined by finitely many polynomial equations.
\end{theorem}

The very first example already displays the algebraicity
\cite[pp.~249--250]{HubbardLundell2011}.

\begin{example}[{\cite[Example~4.2, pp.~249--250]{HubbardLundell2011}}]\label{ex:diffalg:exp}
Let $L(u) = u' - u$ over $\mathbb{C}(t)$. A differential automorphism of
its splitting field sends $e^t$ to $c\,e^t$ for some
$c\in\mathbb{C}^{\times}$, and $c$ determines it. The differential Galois
group is $\mathrm{GL}_1(\mathbb{C})$, the multiplicative group of the
complex numbers.
\end{example}

This is the case drawn in Figure~\ref{fig:diff:pv}.

The heart of the theory is the Galois correspondence
\cite[pp.~29--33]{Kolchin1948PicardVessiot}. Kolchin first shows that a
Picard--Vessiot extension is normal in his sense
\cite[Theorem~1, p.~32]{Kolchin1948PicardVessiot}. The intermediate
differential fields then determine the Zariski-closed subgroups
\cite[Theorem, p.~32]{Kolchin1948PicardVessiot}.

\begin{theorem}[{\cite[$\S$19, Theorem, p.~32]{Kolchin1948PicardVessiot}}]\label{thm:diffalg:pv-correspondence}
Let $G$ be a Picard--Vessiot extension of $F$. The map sending an
intermediate differential field $F_1$ to the group of differential
automorphisms of $G$ over $F_1$ is a bijection between the intermediate
differential fields of $G/F$ and the algebraic subgroups of the Galois
group.
\end{theorem}

Cassidy and Singer package the full result, including the behaviour of
normal subgroups, in one statement
\cite[Theorem~2.1, pp.~5--6]{CassidySinger2007}.

\begin{theorem}[{\cite[Theorem~2.1, pp.~5--6]{CassidySinger2007}}]\label{thm:diffalg:pv-fundamental}
Let $k$ be a $\Delta$-field with algebraically closed field of constants
$C$, and let an integrable system of linear differential equations over
$k$ be given.

\begin{enumerate}
\item There exists a Picard--Vessiot extension $K$ of $k$ associated with
  the system, and this extension is unique up to differential
  $k$-isomorphism.
\item The $\Delta$-Galois group $\mathrm{Gal}_\Delta(K/k)$ equals
  $G(C)$, where $G$ is a linear algebraic group defined over $C$.
\item The map that sends a differential subfield $F$ with
  $k\subset F\subset K$ to $\mathrm{Gal}_\Delta(K/F)$ is a bijection
  between the differential subfields of $K$ containing $k$ and the
  algebraic subgroups of $\mathrm{Gal}_\Delta(K/k)$. Its inverse sends a
  Zariski-closed group $H$ to the fixed field
  $K^H = \{z\in K \mid \sigma(z) = z \text{ for all }\sigma\in H\}$.
\item A Zariski-closed subgroup $H$ of $\mathrm{Gal}_\Delta(K/k)$ is
  normal if and only if the field $K^H$ is invariant under
  $\mathrm{Gal}_\Delta(K/k)$ as a set. Then $K^H$ is a Picard--Vessiot
  extension of $k$ with group $\mathrm{Gal}_\Delta(K/k)/H$.
\end{enumerate}
\end{theorem}

The gain for solvability questions is immediate
\cite[pp.~5--6]{Kolchin1948PicardVessiot}. Kolchin calls the
differential analogue of the radical extension a \emph{Liouvillian}
extension: one built by adjunction of integrals, exponentials of
integrals, and algebraic functions, with no new constant
\cite[pp.~5--6]{Kolchin1948PicardVessiot}.

\begin{theorem}[{\cite[Chapter~V, p.~5]{Kolchin1948PicardVessiot}}]\label{thm:diffalg:pv-liouvillian}
A Picard--Vessiot extension is contained in a Liouvillian extension of
the base differential field if and only if the component of the identity
of its Galois group is solvable.
\end{theorem}

This is the analogue of the criterion of Galois: the Galois group of the
equation decides whether its solutions can be expressed in closed form
\cite[pp.~5--6]{Kolchin1948PicardVessiot}. In the multiplicative example
of Figure~\ref{fig:diff:pv} the group is abelian, hence solvable, and the
solution $e^t$ is literally obtained by exponentiating an integral
\cite[pp.~249--250]{HubbardLundell2011}. In general the criterion of
Theorem~\ref{thm:diffalg:pv-liouvillian} separates the elementary from
the non-elementary.

The same theory runs with an iterative derivation in place of a derivation.
An iterative derivation on a ring $R$ is a sequence of additive maps
$\delta^{(n)}$ \cite[p.~1]{MichelSrinivasan2026}. The first term
$\delta^{(0)}$ is the identity, the Leibniz rule holds in every degree, and
the terms satisfy the binomial rule
$\delta^{(n)}\circ\delta^{(m)}=\binom{m+n}{n}\delta^{(m+n)}$
\cite[p.~1]{MichelSrinivasan2026}. In characteristic zero such a sequence is
determined by its first term \cite[p.~2]{MichelSrinivasan2026}. Michel and
Srinivasan ask when a derivation of the base field extends to a central
simple algebra \cite[p.~1]{MichelSrinivasan2026}.

\begin{theorem}[{\cite[Theorem~2.2, p.~3]{MichelSrinivasan2026}}]\label{thm:diffalg:iterative-extension}
Let $F$ be a field of characteristic $p>0$ and let $A$ be a central simple
$F$-algebra whose exponent in the Brauer group is not divisible by $p$.
Then every iterative derivation on $F$ extends to an iterative derivation
on $A$.
\end{theorem}

The extension fails in general without the hypothesis on the exponent
\cite[p.~2]{MichelSrinivasan2026}. With the extension in hand the authors
build the Picard--Vessiot theory for the algebra: a splitting field exists
and is unique up to isomorphism \cite[p.~1]{MichelSrinivasan2026}. The
Galois group then describes the algebra through its right $\delta$-ideals
\cite[p.~1]{MichelSrinivasan2026}. So the Galois correspondence of this
section survives into the noncommutative coefficient ring.

Section~\ref{sec:differential-algebra:model-theory}
gives a different reading of the same algebra, with model-theoretic
tools.

\hbfigure{17-differential-algebra/galois-correspondence}{The Galois correspondence of Kolchin \cite[$\S$19, Theorem, p.~32]{Kolchin1948PicardVessiot}, in the form of Cassidy and Singer \cite[Theorem~2.1, pp.~5--6]{CassidySinger2007}. An intermediate differential field $F_1$ of the Picard--Vessiot extension $G$ of $F$ corresponds to the algebraic subgroup $\operatorname{Gal}(G/F_1)$, and a subgroup $H$ to its fixed field $G^{H}$. The two maps are inverse bijections and run in opposite directions.}{fig:diff:galois-correspondence}
